%%% Computers in Education Journal Template
%%% Based on Physiome Journal Template 
%%% Creative Commons CC BY 4.0
%%% overleaf.com/latex/templates/physiome-journal-article-template/kfbqwxxmtsfv

%%% Last Modified by CoED Journal Dec 2020


\documentclass[lineno,10pt]{aseecoed}

%% Choose from Original,Letter, Review, Case Study, Methods
\articletype{Original}

%% Main Title
\title{Example Article Title}

%% During Initial Submission please leave the following blank
\author[1][corr.author@xyz.ac.mm]{First Author}
\author[2]{Second Author}
\affil[1]{Address of first author}
\affil[2]{Address of second author}


%% The following lines can be omitted when submitting;
%% information will be added by editors
\publicationdate{30 Sept 2020}
\volume{10}
\issue{4}
\submitteddate{4 May 2020}
\accepteddate{15 Sept 2020}
\citethisas{First and Second (2020) Example Article Title. aseecoed 10(4).}{}

%%% Start of Main Body of Article
\begin{document}

\maketitle

\begin{abstract}
Please provide an abstract of no more than 350 words. Your abstract should explain the main contributions of your article, and should not contain any material that is not included in the main text. 
\end{abstract}

\keywords{Keyword1, Keyword2, Keyword3}

\relatedpubs{sample}{borowczak2018enabling,burrows2017hardening}

\section{Introduction}

Thanks for using Overleaf to write your article. Your introduction goes here! Some examples of commonly used commands and features are listed below, to help you get started.

\subsection{Primary Publication}

Though not required many \emph{aseecoed} articles have associations with one or more related ASEE publications. These publications, whether from the authorship list or not can be listed in the .bib file, then insert them after the \verb|\keywords{...}| using the \verb|\relatedpubs| command:

\verb|\relatedpubs{name of .bib file}{BibTeX keys of the publications}|

If you are authoring and compiling this template on your own machine, you will need to run an extra step \verb|bibtex relatedpubs| to generate them in the final PDF. If you wish, you can use \texttt{latexmk}, \texttt{arara} or \texttt{Make} to automate this step.


\section{Literature Review}
% Explain the current state of the research field with key publications cited.
% Highlight controversial and diverging hypotheses when necessary. 



\subsection{Some \LaTeX{} Examples}
\label{sec:examples}

Use section and subsection commands to organize your document. \LaTeX{} handles all the formatting and numbering automatically. Use \verb|\autoref| and \verb|\label| commands for cross-references, e.g.~\autoref{sec:examples}, \autoref{eq:sum}, \autoref{fig:view}, \autoref{tab:widgets}. You can still use the more common \verb|\ref|, but this will only generate the (sub)section/table/figure/equation number: \ref{tab:withnotes}. 

\subsection{Figures and Tables}

Use the table and tabular commands for basic tables --- see \autoref{tab:widgets}, for example. \autoref{tab:withnotes} shows a larger example with \emph{table notes}. You can upload a figure (JPEG, PNG or PDF) using the project menu. To include it in your document, use the \verb|\includegraphics| command as in the code for \autoref{fig:view} below. Captions are always justified and start from the left; don't try to change the alignment.

If you prefer, you can place all your image files in a folder. Remember to include the folder path in your \verb|\includegraphics| command, or use `\verb|\graphicspath|` to specify the path to the folder in which all your image files can be found.

\begin{figure}[ht]\centering
\includegraphics[width=0.5\linewidth]{images/image1040.png}
\caption{An example image of a logo.}
\label{fig:view}
\end{figure}

\begin{table}[ht]\centering
\caption{An example table.}\label{tab:widgets}
\begin{tabular}{l r}
\toprule
Item & Quantity \\\midrule
Candles & 4 \\
Fork handles & ?\\
\bottomrule
\end{tabular}
\end{table}

\begin{table}[hbt!]\centering
\begin{threeparttable}
\caption{An example table with tablenotes}\label{tab:withnotes}

\begin{tabular}{lccrr}
\toprule
Course & TSC($n$) & Control ($n=40$) & TP & $t$ (68)\\
\midrule
Computer Science I & 38 & 58\tnote{1} & 504.48 & 58 ms\\
Materials and Mechanics & 38 & 58 & 504.48 & 58 ms\\
Differential Equations & 38 & 58 & 504.48 & 58 ms\\
Cybersecurity & 38 & 58 & 504.48 & 58 ms\\
Electronics II\tnote{2} & 38 & 58 & 504.48 & 58 ms\\
Basket weaving & 38 & 58 & 504.48 & 58 ms\\
\bottomrule
\end{tabular}
\begin{tablenotes}
\item[1] here's a note.
\item[2] and another.
\end{tablenotes}
\end{threeparttable}
\end{table}

\subsection{Citations}

LaTeX formats citations and references automatically using the bibliography records in your .bib file, which you can edit via the project menu. Use the \verb|\citet| command for a text citation, like \citet{borowczak2019mitigating}, and the \verb|\citep| command for a citation in parentheses \citep{burrows2018integrated}.

\subsection{Mathematics}

\LaTeX{} is great at typesetting mathematics. Let $X_1, X_2, \ldots, X_n$ be a sequence of independent and identically distributed random variables with $\text{E}[X_i] = \mu$ and $\text{Var}[X_i] = \sigma^2 < \infty$, and let
\begin{equation}\label{eq:sum}
S_n = \frac{X_1 + X_2 + \cdots + X_n}{n}
      = \frac{1}{n}\sum_{i}^{n} X_i
\end{equation}
denote their mean. Then as $n$ approaches infinity, the random variables $\sqrt{n}(S_n - \mu)$ converge in distribution to a normal $\mathcal{N}(0, \sigma^2)$.

\subsection{Lists}

You can make lists with automatic numbering \dots

\begin{enumerate}[noitemsep] 
\item Like this,
\item and like this.
\end{enumerate}
\dots or bullet points \dots
\begin{itemize}[noitemsep] 
\item Like this,
\item and like this.
\end{itemize}
\dots or with words and descriptions \dots
\begin{description}
\item[Word] Definition
\item[Concept] Explanation
\item[Idea] Text
\end{description}

\section{Methods and Context}
% Methods and protocols should be described in detail and well-established methods can be briefly described and appropriately cited. Provide details that enable readers to frame/situate this work - locations, participants demographics, details relevant to your particular study/implementation
 Methods and protocols should be described in detail and well-established methods can be briefly described and appropriately cited. Provide details that enable readers to frame/situate this work - locations, participants demographics, details relevant to your particular study/implementation

Quam suscipit ut quidem et animi numquam consectetur et. Nihil et commodi ut officia eveniet beatae qui. Placeat accusantium eius consequatur animi nisi sed. Pariatur et dolores tempore velit similique voluptatem similique error.

Quam suscipit ut quidem et animi numquam consectetur et. Nihil et commodi ut officia eveniet beatae qui. Placeat accusantium eius consequatur animi nisi sed. Pariatur et dolores tempore velit similique voluptatem similique error. Quam suscipit ut quidem et animi numquam consectetur et. Nihil et commodi ut officia eveniet beatae qui. Placeat accusantium eius consequatur animi nisi sed. Pariatur et dolores tempore velit similique voluptatem similique error.

\section{Discussion}
%Authors should discuss the results and how they can be interpreted in perspective of previous studies and of the working hypotheses. 

Authors should discuss the results and how they can be interpreted in perspective of previous studies and of the working hypotheses. 

 Quam suscipit ut quidem et animi numquam consectetur et. Nihil et commodi ut officia eveniet beatae qui. Placeat accusantium eius consequatur animi nisi sed. Pariatur et dolores tempore velit similique voluptatem similique error.

 Quam suscipit ut quidem et animi numquam consectetur et. Nihil et commodi ut officia eveniet beatae qui. Placeat accusantium eius consequatur animi nisi sed. Pariatur et dolores tempore velit similique voluptatem similique error.

\section{Conclusion}
%Be sure to include the implications of your work that deals with computing for (Engineering) Educators!

Be sure to include the implications of your work that deals with computing for (Engineering) Educators!


\bibliography{sample}

\end{document}